\scanpage{111}
\deleted{\unsure{$H$},} Un $H\inserted{\supset1}$ est défini aussi par son orthogonal, et de $=$ pour $H'$
\[
V\subset A^{\unsure{s|\leftrightarrow|\natural}},\quad V'\subset A^s
\]
sous fibrés vectoriels de rang $2$, qui satisfont donc
\[
V+V'=A^s
\]
donc
\[
V\cap V'\overset{\text{def}}{=}D\qquad\unsure{\text{où}}
\]
est un sous fibré vect. de rang $1$, $\unsure{*}$

\deleted{Dire $H$ est étale} \deleted{que $H$ est}

Les $H$, correspondent à des sections \inserted{sur $k$ ou $k$-points} du \deleted{fibré} \inserted{plan} projectif \inserted{\unsure{duel|dual}} \deleted{\unsure{$V^\perp$}} $\check{\mathbb{P}}(A/k)$, et les $V$ à des droites dans le plan projectif dual, dans lequel il y a une conique canonique, celle des droites isotropes de $A^s$. Alors $H$ est \unsure{étale|étal|ét.} sur $k$ (\unsure{en tt que|les $H$ qui} algèbre sur $k$) sss en chaque \deleted{point} fibre \inserted{géom.}, la droite correspondante (polaire) $D_H=D_V$ est tangente à la conique, i.e. y découpe un diviseur \deleted{de} non réduit : un pt. Si cette condition est vérifiée pour
% slanted
\margin{\deleted{\uline{NB} Les $H$ en question sont en corr. $1$-$1$ avec les diviseurs de degré $2$ sur la conique qui sont non dégénérés.}}
% archive sheet number: 110
